% Short JSON Lab excerpt: the lab-level wrapper plus the head of module 1.
% Source: Artifacts/Data/Benchmark/Outputs/
%         claude-sonnet-5_json_lab_L3_20260710_154838_output.json
%         (lab = vectors_and_projectile_motion, code level L3 = paper "Objectives",
%          structure = multi-module, 4 modules).
% Numeric arrays collapsed onto one line; elisions marked with a bare ellipsis.
% Content is otherwise verbatim.
%
% Requires: \usepackage{listings} \usepackage{xcolor} \usepackage{textcomp}
%           and the arlemlisting style defined in lab_topics_appendix.tex.

\begin{lstlisting}[style=arlemlisting, caption={The JSON Lab wrapper and the opening of the first module, from Claude Sonnet~5's \emph{Vectors and Projectile Motion} lab. The lab object carries course-level metadata and objectives; each module declares its own scene objects and the clips that animate them.}, label={lst:json-lab-snippet}]
{
  "version": "2.0",
  "labId": "phys101_vectors_projectile_motion",
  "author": "A Robot",
  "institution": "MTSU",
  "courseName": "General Physics I (Mechanics)",
  "estimatedLength": "40 minutes",
  "objectives": [
    "Represent 2D vectors by magnitude/direction and by horizontal/vertical components.",
    "Add vectors graphically using the tip-to-tail method to find a resultant.",
    "Decompose vectors into components using right-triangle trigonometry and reconstruct vectors from components.",
    "Explain the independence of horizontal and vertical motion in projectile motion.",
    "Describe how velocity components and position change throughout a projectile's flight.",
    "Predict how launch angle and speed affect trajectory, range, and maximum height."
  ],
  "modules": [
    {
      "moduleType": "demo",
      "prefab": "demoPrefab",
      "moduleName": "Vector Basics: Magnitude, Direction, and Addition",
      "description": "Introduces vectors as arrows with magnitude and direction, and demonstrates tip-to-tail vector addition to find a resultant.",
      "author": "A Robot",
      "institution": "MTSU",
      "dateCreated": "2026-07-10",
      "educationalObjectives": [
        "Represent a 2D vector by its magnitude and direction.",
        "Add two vectors graphically using the tip-to-tail method.",
        "Identify the resultant vector of a vector sum."
      ],
      "instructions": [
        "Observe the two vectors labeled A and B.",
        "Watch as vector B is moved tip-to-tail with vector A to form the resultant vector R."
      ],
      "objects": [
        {
          "prefab": "vectorArrowPrefab",
          "name": "vectorA",
          "position": [-0.5, 0, 0.3],
          "scale": [0.3, 0.02, 0.02],
          "color": [220, 40, 40, 255],
          "components": [
            {
              "componentType": "textMeshPro",
              "text": "A",
              "color": [220, 40, 40, 255],
              "fontSize": 28
            }
          ]
        },
        ...
      ],
      "clips": [
        {
          "clipName": "clip0",
          "audioClip": "vectors_intro",
          "autoAdvance": true,
          "changeMeaning": "Initial staging: two vectors A and B appear at separate positions with labels, and a caption introduces the concept of vectors.",
          "narration": "In physics, a vector is a quantity that has both magnitude and direction, like velocity or force. Here we see two vectors, A and B, represented as arrows. The length of each arrow shows its magnitude, and the way it points shows its direction."
        },
        ...
      ]
    },
    ...
  ]
}
\end{lstlisting}
